\documentclass[11pt,a4paper]{article}

\usepackage[utf8]{inputenc}
\usepackage[T1]{fontenc}
\usepackage[french]{babel}
\usepackage{amsmath,amssymb,amsfonts}
\usepackage{geometry}
\usepackage{graphicx}
\usepackage{booktabs}
\usepackage{hyperref}
\usepackage{array}
\usepackage{xcolor}
\usepackage{float}

\geometry{margin=2.5cm}
\graphicspath{{figures/}}

\title{
\textbf{Corrections, prédictions et déviations observables}\\
\large Le tenseur \(\mathcal{H}_{\mu\nu}\) dans Bottom-Up Quantum Gravity
}

\author{Farid Hamdad}
\date{\today}

\begin{document}

\maketitle

\begin{abstract}
Paper~20 étudie les corrections à la limite d'Einstein effective de
Bottom-Up Quantum Gravity. Paper~19 a proposé l'équation
\[
G_{\mu\nu}[g^{\rm ent}]
+
\Lambda_{\rm ent}g_{\mu\nu}^{\rm ent}
=
8\pi G_{\rm eff}T_{\mu\nu}^{\rm ent}
+
\mathcal{H}_{\mu\nu}.
\]
Le présent papier analyse le tenseur de correction
\(\mathcal{H}_{\mu\nu}\). Nous le décomposons en secteurs spectral, courbure,
source, dimensionnel, non local, topologique et taille finie :
\[
\mathcal{H}_{\mu\nu}
=
\mathcal{H}^{\rm spec}_{\mu\nu}
+
\mathcal{H}^{\rm curv}_{\mu\nu}
+
\mathcal{H}^{\rm source}_{\mu\nu}
+
\mathcal{H}^{\rm dim}_{\mu\nu}
+
\mathcal{H}^{\rm nonlocal}_{\mu\nu}
+
\mathcal{H}^{\rm topo}_{\mu\nu}
+
\mathcal{H}^{\rm finite}_{\mu\nu}.
\]
Les secteurs spectral, courbure et source sont contraints par Papers~16--18.
Les secteurs dimensionnel et non local sont reliés aux résultats
phénoménologiques SPARC et lensing. Le secteur dimensionnel est associé à une
fenêtre effective \(d_{\rm eff}\sim2.02\to2.64\), tandis que le secteur non
local est relié à une amplification lensing moyenne
\(\langle{\rm ratio}_{\rm peak}\rangle=2.721481\). Paper~20 identifie ainsi
les principales déviations observables de BuP par rapport à la limite
d'Einstein lisse.
\end{abstract}

\tableofcontents

\section{Introduction}

Paper~19 a assemblé les trois briques nécessaires à la limite d'Einstein de
Bottom-Up Quantum Gravity :
\[
L_N\to\Delta_g,
\qquad
\kappa^{OR}\to R_{\mu\nu}u^\mu u^\nu,
\qquad
\delta W_{\rm loc}\to T_{\mu\nu}^{\rm ent}.
\]

L'équation effective obtenue est :
\[
G_{\mu\nu}[g^{\rm ent}]
+
\Lambda_{\rm ent}g_{\mu\nu}^{\rm ent}
=
8\pi G_{\rm eff}T_{\mu\nu}^{\rm ent}
+
\mathcal{H}_{\mu\nu}.
\]

Dans la limite lisse, locale et basse énergie :
\[
\mathcal{H}_{\mu\nu}\to0.
\]

BuP se réduit alors à :
\[
G_{\mu\nu}
+
\Lambda_{\rm ent}g_{\mu\nu}
=
8\pi G_{\rm eff}T_{\mu\nu}^{\rm ent}.
\]

Mais dans les régimes où l'intrication n'est pas parfaitement lisse, locale ou
homogène, le tenseur
\[
\mathcal{H}_{\mu\nu}
\]
devient non nul. C'est ce tenseur qui porte les prédictions différenciant BuP
de la relativité générale standard et de \(\Lambda\)CDM.

\section{Décomposition du tenseur de correction}

Nous écrivons :
\[
\mathcal{H}_{\mu\nu}
=
\mathcal{H}_{\mu\nu}^{\rm spec}
+
\mathcal{H}_{\mu\nu}^{\rm curv}
+
\mathcal{H}_{\mu\nu}^{\rm source}
+
\mathcal{H}_{\mu\nu}^{\rm dim}
+
\mathcal{H}_{\mu\nu}^{\rm nonlocal}
+
\mathcal{H}_{\mu\nu}^{\rm topo}
+
\mathcal{H}_{\mu\nu}^{\rm finite}.
\]

Chaque terme correspond à une origine physique différente :

\[
\begin{array}{lll}
\toprule
{\rm Secteur} & {\rm Origine} & {\rm Cible}\\
\midrule
\mathcal{H}^{\rm spec}_{\mu\nu} & {\rm spectre\ de\ }L & {\rm corrections\ de\ courbure\ supérieure}\\
\mathcal{H}^{\rm curv}_{\mu\nu} & \kappa^{OR}\to Ricci & {\rm calibration\ de\ courbure}\\
\mathcal{H}^{\rm source}_{\mu\nu} & \delta\langle K_A\rangle & {\rm reconstruction\ de\ source}\\
\mathcal{H}^{\rm dim}_{\mu\nu} & d_s,d_w,\alpha_{\rm eff} & {\rm galaxies,\ cosmologie}\\
\mathcal{H}^{\rm nonlocal}_{\mu\nu} & {\rm liens\ longs\ d'intrication} & {\rm lensing,\ SLACS}\\
\mathcal{H}^{\rm topo}_{\mu\nu} & S_{\rm topo}[W] & {\rm holonomies,\ topologie}\\
\mathcal{H}^{\rm finite}_{\mu\nu} & N<\infty & {\rm scaling\ numérique}\\
\bottomrule
\end{array}
\]

\section{Secteur spectral}

Le secteur spectral provient du terme :
\[
\mathrm{Tr}\,L^{-\beta}.
\]

Par représentation de Mellin :
\[
\mathrm{Tr}\,L^{-\beta}
=
\frac{1}{\Gamma(\beta)}
\int_0^\infty
t^{\beta-1}
\mathrm{Tr}(e^{-tL})\,dt.
\]

Dans la limite continue :
\[
L\to-\Delta_g.
\]

La trace de chaleur s'écrit :
\[
\mathrm{Tr}(e^{-t\Delta_g})
\sim
(4\pi t)^{-D/2}
\left(
a_0+a_1t+a_2t^2+\cdots
\right).
\]

Le terme \(a_1\) contient la courbure scalaire et conduit au secteur
Einstein--Hilbert. Les termes suivants génèrent :
\[
R^2,
\qquad
R_{\mu\nu}R^{\mu\nu},
\qquad
R_{\mu\nu\rho\sigma}R^{\mu\nu\rho\sigma},
\qquad
\nabla^2R,
\qquad
\cdots.
\]

Le budget de corrections donne :
\[
\langle |H_{\rm spec}|\rangle=0.015850,
\]
avec :
\[
\max |H_{\rm spec}|=0.026627.
\]

Le secteur spectral est donc borné numériquement à quelques pourcents sur les
géométries contrôlées.

\section{Secteur courbure}

Paper~17 a établi une calibration moyenne :
\[
\frac{\kappa^{OR}}{\epsilon}
\simeq
B_N+C_N R_{\mu\nu}u^\mu u^\nu.
\]

À \(N=512\), on obtient :
\[
B_N=-0.301281,
\]
\[
C_N=0.391933,
\]
\[
A_N=\frac{1}{C_N}=2.551455.
\]

La correction de courbure est donc :
\[
\mathcal{H}_{\mu\nu}^{\rm curv}
\sim
{\rm biais\ fini}
+
{\rm renormalisation\ OR\to Ricci}.
\]

Ce secteur n'est pas une anomalie. Il exprime le fait que la courbure
d'Ollivier--Ricci discrète nécessite une calibration pour se convertir en
courbure de Ricci continue.

\section{Secteur source}

Paper~18 a montré :
\[
\delta S_A^{\rm graph}
\simeq
\delta\langle K_A^{\rm graph}\rangle.
\]

Sur le tore plat :
\[
\delta S_A\simeq0.989481\,\delta\langle K_A\rangle,
\]
et sur la sphère :
\[
\delta S_A\simeq0.989718\,\delta\langle K_A\rangle.
\]

Le résidu moyen du secteur modulaire est donc :
\[
H_{\rm source\_modular}\simeq0.010400.
\]

Paper~18 a aussi montré :
\[
|\delta\langle K_A\rangle|
\to
\langle|\Delta\kappa|\rangle_{\rm near}.
\]

Le résidu source--courbure est :
\[
H_{\rm source\_curvature}\simeq0.025260.
\]

Ainsi, le secteur source est bien contrôlé :
\[
\mathcal{H}_{\mu\nu}^{\rm source}
\]
représente un petit écart entre la variation modulaire idéale et la réponse de
courbure mesurée.

\section{Secteur dimensionnel et SPARC}

Le secteur dimensionnel est associé aux variations de dimension effective :
\[
d_s(x),
\qquad
d_w(x),
\qquad
\alpha_{\rm eff}(x).
\]

Il peut s'écrire formellement :
\[
\mathcal{H}_{\mu\nu}^{\rm dim}
=
\mathcal{H}_{\mu\nu}
[
d_s(x),
d_w(x),
\alpha_{\rm eff}(x)
].
\]

Des termes typiques peuvent prendre la forme :
\[
\nabla_\mu\nabla_\nu d_s
-
g_{\mu\nu}\Box d_s,
\]
ou :
\[
(\nabla_\mu d_s)(\nabla_\nu d_s).
\]

Le budget v3 relie ce secteur aux courbes de rotation SPARC. Le résultat
global est :
\[
{\rm median}\ \chi^2_{\rm red}=12.568829,
\]
\[
{\rm fraction}(\chi^2_{\rm red}<2)=0.182857,
\]
\[
{\rm fraction}(\chi^2_{\rm red}<10)=0.485714.
\]

La taxonomie de qualité donne :
\[
f_{\rm excellent}=0.182857,
\]
\[
f_{\rm good}=0.154286,
\]
\[
f_{\rm medium}=0.274286,
\]
\[
f_{\rm poor}=0.388571.
\]

Donc :
\[
f_{\rm excellent+good+medium}=0.611429.
\]

La fenêtre dimensionnelle mesurée est :
\[
\langle d_{\rm used,min}\rangle=2.023651,
\]
\[
\langle d_{\rm used,max}\rangle=2.643653.
\]

Ainsi, \(61.1\%\) des galaxies sont au moins dans un régime exploitable jusqu'à
medium, tandis que \(38.9\%\) restent dans un régime pauvre. Ces galaxies
pauvres ne doivent pas être ignorées : elles signalent précisément les
corrections morphologiques, de phase ou non locales encore absentes.

\section{Secteur non local et lentilles}

Le secteur non local vient des liens longs du graphe d'intrication :
\[
W_{ij}\neq0
\]
pour des nœuds éloignés dans la géométrie émergente.

Dans la limite continue, il peut prendre une forme intégrale :
\[
\mathcal{H}_{\mu\nu}^{\rm nonlocal}(x)
\sim
\int
K_{\mu\nu}^{\ \ \alpha\beta}(x,y)
T_{\alpha\beta}(y)
d^Dy.
\]

Ce secteur peut produire des effets de type matière noire effective :
\[
{\rm rotation\ curves},
\qquad
{\rm lensing},
\qquad
{\rm amplification\ non\ locale}.
\]

Le budget v3 relie ce secteur au bridge lensing. Les résultats sont :
\[
\langle{\rm ratio}_{\rm peak}\rangle=2.721481,
\]
\[
{\rm median}({\rm ratio}_{\rm peak})=1.967831,
\]
\[
\max({\rm ratio}_{\rm peak})=13.807607.
\]

L'amplitude absolue de réponse est :
\[
\langle\Delta_{\rm peak}\rangle=889.087288,
\]
\[
{\rm median}(\Delta_{\rm peak})=704.916728,
\]
\[
\max(\Delta_{\rm peak})=3594.219422.
\]

La fraction de profils lensing valides est :
\[
f_{\rm ok}=0.977143.
\]

Ainsi :
\[
\mathcal{H}_{\mu\nu}^{\rm nonlocal}
\]
est maintenant relié quantitativement à la phénoménologie lensing.

\section{Branche lensing historique et SLACS}

Avant la formulation du budget
\[
\mathcal{H}_{\mu\nu},
\]
BuP contenait déjà une branche lensing, notamment motivée par les tests de
type SLACS et par l'étude de spirales massives.

Cette branche suivait la chaîne :
\[
d(r)
\to
f_d(r)
\to
g_{\rm eff}(r)
\to
\alpha(R).
\]

Le facteur dimensionnel était :
\[
f_d(r)
=
\left(
\frac{r}{r_{\rm ref}}
\right)^{3-d(r)}.
\]

La réponse optique effective était :
\[
g_{\rm eff}(r)
=
g_{\rm bary}(r)
\left[
1+
\lambda w_{\rm tot}(r)
(f_d(r)-1)
\right].
\]

Un régime optique centralisé avait été identifié :
\[
(\lambda,\mathrm{width},x_{\rm cap},r_{\rm cut})
=
(3.0,0.8,4.0,5.0).
\]

Ce régime produisait des excès robustes de déflexion typiquement de l'ordre
de \(30\%\) à \(50\%\) sur plusieurs spirales massives.

Dans le langage de Paper~20, cette ancienne branche se réinterprète comme :
\[
\mathcal{H}_{\mu\nu}^{\rm nonlocal/optical}.
\]

Ainsi :
\[
\boxed{
{\rm branche\ lensing/SLACS}
\subset
\mathcal{H}_{\mu\nu}^{\rm nonlocal}.
}
\]

Le budget v3 ne remplace donc pas l'ancienne branche lensing : il la formalise
comme un secteur de correction non locale.

\section{Secteurs topologique et taille finie}

Deux secteurs restent encore partiellement ouverts.

Le secteur topologique :
\[
\mathcal{H}_{\mu\nu}^{\rm topo}
\]
doit être relié à :
\[
S_{\rm topo}[W],
\]
aux holonomies, aux constantes modulaires et aux contraintes globales du
graphe d'intrication.

Le secteur de taille finie :
\[
\mathcal{H}_{\mu\nu}^{\rm finite}
\]
doit être relié aux lois de scaling :
\[
N\to\infty,
\]
\[
\rho_{\rm graph}\to\rho_{\rm continuum}.
\]

Ces deux secteurs seront importants pour distinguer les effets physiques des
artefacts numériques.

\section{Résumé du budget v3}

\[
\begin{array}{llll}
\toprule
{\rm Secteur} & {\rm Moyenne} & {\rm Maximum} & {\rm Statut}\\
\midrule
H_{\rm spec} & 0.015850 & 0.026627 & {\rm borné}\\
H_{\rm source\_modular} & 0.010400 & 0.010519 & {\rm petit\ résidu}\\
H_{\rm source\_curvature} & 0.025260 & 0.032771 & {\rm petit\ résidu}\\
H_{\rm dim} & 8.508743 & 12.568829 & {\rm SPARC}\\
H_{\rm nonlocal} & 445.904384 & 889.087288 & {\rm lensing}\\
H_{\rm localization} & 7.913268 & 9.533529 & {\rm localité\ positive}\\
\bottomrule
\end{array}
\]

Les valeurs associées à \(H_{\rm dim}\) et \(H_{\rm nonlocal}\) doivent être
lues comme des proxys phénoménologiques et non comme des normes tensorielles.

\section{Interprétation physique}

Paper~20 montre que les déviations BuP ne sont pas arbitraires. Elles sont
organisées en secteurs :

\[
\mathcal{H}_{\mu\nu}^{\rm spec}
\]
contrôle les corrections géométriques de haute énergie ;

\[
\mathcal{H}_{\mu\nu}^{\rm dim}
\]
contrôle les écarts galactiques et cosmologiques liés à la dimension
effective ;

\[
\mathcal{H}_{\mu\nu}^{\rm nonlocal}
\]
contrôle les amplifications de type lensing et les effets de liens longs ;

\[
\mathcal{H}_{\mu\nu}^{\rm source}
\]
contrôle les écarts de reconstruction source--courbure ;

\[
\mathcal{H}_{\mu\nu}^{\rm topo}
\]
encode les contraintes globales ;

\[
\mathcal{H}_{\mu\nu}^{\rm finite}
\]
encode les artefacts de discrétisation.

Ainsi, BuP ne prédit pas seulement une limite d'Einstein. Il prédit aussi une
structure de corrections observables.

\section{Conclusion}

Paper~20 classe les corrections à la limite d'Einstein effective de BuP.

Le résultat principal est :
\[
\mathcal{H}_{\mu\nu}
=
\mathcal{H}^{\rm spec}_{\mu\nu}
+
\mathcal{H}^{\rm curv}_{\mu\nu}
+
\mathcal{H}^{\rm source}_{\mu\nu}
+
\mathcal{H}^{\rm dim}_{\mu\nu}
+
\mathcal{H}^{\rm nonlocal}_{\mu\nu}
+
\mathcal{H}^{\rm topo}_{\mu\nu}
+
\mathcal{H}^{\rm finite}_{\mu\nu}.
\]

Les secteurs spectral et source sont déjà petits et bien contrôlés. Le secteur
dimensionnel est relié à SPARC et à une fenêtre effective
\[
d_{\rm eff}\sim2.02\to2.64.
\]
Le secteur non local est relié aux lentilles, avec une amplification moyenne
\[
\langle{\rm ratio}_{\rm peak}\rangle=2.721481.
\]

Paper~20 donne donc la première carte structurée des déviations observables de
BuP par rapport à la limite d'Einstein lisse.

\section*{Disponibilité du code et des données}

Le dossier principal est :
\[
\texttt{papers/paper20\_corrections\_predictions/}.
\]

Les scripts principaux sont :
\[
\texttt{scripts/paper20\_build\_correction\_budget\_v1.py},
\]
\[
\texttt{scripts/paper20\_build\_correction\_budget\_v2.py},
\]
\[
\texttt{scripts/paper20\_build\_correction\_budget\_v3.py}.
\]

Les résultats principaux sont dans :
\[
\texttt{results/correction\_budget\_v3\_sparc\_lensing/}.
\]

\appendix

\section{Prochaines étapes}

Les prochaines étapes sont :

\begin{enumerate}
    \item relier \(H_{\rm topo}\) aux holonomies et constantes modulaires ;
    \item relier \(H_{\rm finite}\) aux lois de scaling en \(N\) ;
    \item intégrer les résultats SLACS dans un fichier standardisé ;
    \item séparer \(H_{\rm dim}\) par morphologie galactique ;
    \item séparer \(H_{\rm nonlocal}\) par type de lentille ;
    \item relier les corrections cosmologiques à \(H_{\rm dim}+H_{\rm IR}\) ;
    \item produire une version publication-grade du budget de corrections.
\end{enumerate}

\end{document}
